% =============================================================
% 子文件模板：题目描述部分
% 用法：在主文件中使用 \input{filename.tex} 导入
% =============================================================

\begin{center}
    % 建议格式：序号. 中文题目名（英文题目名）
    \section{B. 分糖果 (candy)}
\end{center}

\subsection{题目描述}

% 在此处填写题目背景和详细描述
% 行内公式请使用 $...$，例如 $a+b$
云小斗要给 $n$ 个小朋友分 $m$ 个糖果。他希望能平分，因此保证 $m$ 能被 $n$ 整除，记平均数为 $A=\dfrac{m}{n}$。

但每个小朋友的口味不同：第 $i$ 个小朋友希望分到的糖果数在 $[l_i, r_i]$ 之间。
设第 $i$ 个小朋友实际分到 $x_i$ 个糖果，则不满意值定义为 $\sum_{i=1}^{n} |A-x_i|$。**注意，不满意值是对整个分糖序列的不满意值而不是对每个小朋友。**
请你求出 **最小可能的不满意值**。如果不存在满足所有区间限制且总和为 $m$ 的分配方案，则输出 **-1**。

\subsection{输入格式}

% 描述输入数据的格式，例如：第一行包含一个整数 $n$。
第一行 $n,m$。

后面 $n$ 行，第 $i+1$ 行代表 $l_i,r_i$。

\subsection{输出格式}

% 描述输出数据的格式
一个数，代表最小 $\sum_{i=1}^{n} (\frac{m}{n}-x_i)$。

\subsection{样例输入1}

% 将样例输入直接粘贴在下方，无需额外格式
\begin{lstlisting}
2 10
3 4
5 7
\end{lstlisting}

\subsection{样例输出1}

% 将样例输出直接粘贴在下方
\begin{lstlisting}
2
\end{lstlisting}

\subsection{样例输入2}

% 将样例输入直接粘贴在下方，无需额外格式
\begin{lstlisting}
3 18 
1 3
2 6
3 11
\end{lstlisting}

\subsection{样例输出2}

% 将样例输出直接粘贴在下方
\begin{lstlisting}
6
\end{lstlisting}

\subsection{数据范围和约定}

% 描述数据规模和特殊限制
% 示例：对于 $100\%$ 的数据，$1 \leq n \leq 100$。
对于 $20\%$ 的数据，$1 \le n \le 100,1 \le l_i \le r_i \le 200$。

对于 $40\%$ 的数据，$1 \le n \le 1000$。

对于另外 $10\%$ 的数据，$\frac{m}{n} = 1$。

对于另外 $10\%$ 的数据，保证 $l_i = r_i$。

对于 $100\%$ 的数据，$1 \le n \le 5 \times 10^5,1 \le m \le 10^{15},1 \le \frac{n}{m} \le 10^{10},1 \le l_i \le r_i \le 10^{10}$。
